\chapter{Advanced Guide --- Bijou}\label{record-001}

\noindent{\footnotesize\color{BookMuted}\path{ADVANCED_GUIDE.md}\par}\smallskip

This is the second-track manual for Bijou.

\begin{center}
\BookDiagram{\begin{tikzpicture}[box/.style={draw=BookHeading,fill=BookPanel,rounded corners=1pt,text width=38mm,inner sep=5pt,align=left,font=\footnotesize},edge/.style={-{Latex[length=2mm]},draw=BookHeading}]
\node[box] (S1) at (-5.0,0.0) {Surface Definition / @flyingrobots/bijou (Foundation)};
\node[box] (S2) at (0.0,0.0) {Component Rendering / @flyingrobots/bijou (Foundation)};
\node[box] (T1) at (0.0,-2.6) {Layout Stage / @flyingrobots/bijou-tui (Engine)};
\node[box] (T2) at (0.0,-5.2) {Paint Stage / @flyingrobots/bijou-tui (Engine)};
\node[box] (T3) at (0.0,-7.800000000000001) {Post-Process Shaders / @flyingrobots/bijou-tui (Engine)};
\node[box] (T4) at (0.0,-10.4) {Diffing Layer / @flyingrobots/bijou-tui (Engine)};
\node[box] (N1) at (0.0,-13.0) {ANSI Emission / @flyingrobots/bijou-node (Host)};
\node[box] (N2) at (5.0,0.0) {Worker Orchestration / @flyingrobots/bijou-node (Host)};
\draw[edge,] (S2) to (T1);
\draw[edge,] (T1) to (T2);
\draw[edge,] (T2) to (T3);
\draw[edge,] (T3) to (T4);
\draw[edge,] (T4) to (N1);
\end{tikzpicture}}
\end{center}

Use the root \href{https://github.com/flyingrobots/bijou/blob/cbd70612140e7a450d764b3c7be1c6af878b99bc/README.md}{README} and the package \texttt{G\allowbreak{}U\allowbreak{}I\allowbreak{}D\allowbreak{}E\allowbreak{}.\allowbreak{}m\allowbreak{}d} files when you want the productive-fast path. Use this guide when you need the deeper doctrine, the proving surfaces, or the repo workflows that sit behind the public quick starts.

\section{What Belongs Here}\label{record-001-what-belongs-here}

This guide is for topics like:

\begin{itemize}
\tightlist
\item
  bytes-only internal rendering rules and explicit string boundaries
\item
  shell doctrine and \texttt{c\allowbreak{}r\allowbreak{}e\allowbreak{}a\allowbreak{}t\allowbreak{}e\allowbreak{}F\allowbreak{}r\allowbreak{}a\allowbreak{}m\allowbreak{}e\allowbreak{}d\allowbreak{}A\allowbreak{}p\allowbreak{}p\allowbreak{}(\allowbreak{})} ownership
\item
  programmable pipeline, shaders, paint, transforms, and middleware seams
\item
  DOGFOOD as the canonical proving surface
\item
  performance and soak workflows
\item
  localization workflows and bidirectionality posture
\item
  worker/runtime boundaries and scripted capture
\item
  smoke, PTY, canary, and release-readiness expectations
\end{itemize}

\section{Reading Order}\label{record-001-reading-order}

Use this sequence when you need repo-truth beyond the quick path:

\begin{enumerate}
\def\labelenumi{\arabic{enumi}.}
\tightlist
\item
  Read the \href{https://github.com/flyingrobots/bijou/blob/cbd70612140e7a450d764b3c7be1c6af878b99bc/README.md}{Root README} for the package map and public posture.
\item
  Read the package guide you are actively using.
\item
  Read the matching package \texttt{A\allowbreak{}D\allowbreak{}V\allowbreak{}A\allowbreak{}N\allowbreak{}C\allowbreak{}E\allowbreak{}D\allowbreak{}\_\allowbreak{}G\allowbreak{}U\allowbreak{}I\allowbreak{}D\allowbreak{}E\allowbreak{}.\allowbreak{}m\allowbreak{}d}.
\item
  Use DOGFOOD, the performance docs, and the strategy/spec docs as the proving surfaces for the subsystem you are touching.
\end{enumerate}

\section{Package Lanes}\label{record-001-package-lanes}

\begin{itemize}
\tightlist
\item
  Core advanced work: \href{https://github.com/flyingrobots/bijou/blob/cbd70612140e7a450d764b3c7be1c6af878b99bc/packages/bijou/ADVANCED_GUIDE.md}{packages/bijou/ADVANCED\_GUIDE.md}
\item
  Runtime advanced work: \href{https://github.com/flyingrobots/bijou/blob/cbd70612140e7a450d764b3c7be1c6af878b99bc/packages/bijou-tui/ADVANCED_GUIDE.md}{packages/bijou-tui/ADVANCED\_GUIDE.md}
\item
  Node/runtime-boundary work: \href{https://github.com/flyingrobots/bijou/blob/cbd70612140e7a450d764b3c7be1c6af878b99bc/packages/bijou-node/ADVANCED_GUIDE.md}{packages/bijou-node/ADVANCED\_GUIDE.md}
\item
  Localization runtime: \href{https://github.com/flyingrobots/bijou/blob/cbd70612140e7a450d764b3c7be1c6af878b99bc/packages/bijou-i18n/README.md}{packages/bijou-i18n/README.md}
\item
  Localization tooling: \href{https://github.com/flyingrobots/bijou/blob/cbd70612140e7a450d764b3c7be1c6af878b99bc/packages/bijou-i18n-tools/README.md}{packages/bijou-i18n-tools/README.md}
\item
  Filesystem exchange tooling: \href{https://github.com/flyingrobots/bijou/blob/cbd70612140e7a450d764b3c7be1c6af878b99bc/packages/bijou-i18n-tools-node/README.md}{packages/bijou-i18n-tools-node/README.md}
\item
  Spreadsheet exchange tooling: \href{https://github.com/flyingrobots/bijou/blob/cbd70612140e7a450d764b3c7be1c6af878b99bc/packages/bijou-i18n-tools-xlsx/README.md}{packages/bijou-i18n-tools-xlsx/README.md}
\end{itemize}

\section{Rendering Doctrine}\label{record-001-rendering-doctrine}

The core rendering posture is simple:

\begin{itemize}
\tightlist
\item
  structured terminal UI should stay on the \texttt{S\allowbreak{}u\allowbreak{}r\allowbreak{}f\allowbreak{}a\allowbreak{}c\allowbreak{}e} path
\item
  strings are explicit boundary output, not the internal runtime model
\item
  themes, backgrounds, and shell chrome should resolve to bytes before the hot render loop, not by reparsing strings every frame
\item
  lowering back to text is allowed, but it should happen deliberately at the boundary
\end{itemize}

Start here when that doctrine matters:

\begin{itemize}
\tightlist
\item
  \href{https://github.com/flyingrobots/bijou/blob/cbd70612140e7a450d764b3c7be1c6af878b99bc/packages/bijou/ADVANCED_GUIDE.md}{packages/bijou/ADVANCED\_GUIDE.md}
\item
  \href{https://github.com/flyingrobots/bijou/blob/cbd70612140e7a450d764b3c7be1c6af878b99bc/packages/bijou-tui/ADVANCED_GUIDE.md}{packages/bijou-tui/ADVANCED\_GUIDE.md}
\item
  \href{https://github.com/flyingrobots/bijou/blob/cbd70612140e7a450d764b3c7be1c6af878b99bc/docs/perf/RE-017-byte-pipeline.md}{Byte-Pipeline Recovery}
\item
  \href{https://github.com/flyingrobots/bijou/blob/cbd70612140e7a450d764b3c7be1c6af878b99bc/docs/strategy/low-allocation-renderer.md}{Low-Allocation Renderer}
\item
  \href{https://github.com/flyingrobots/bijou/blob/cbd70612140e7a450d764b3c7be1c6af878b99bc/docs/specs/rendering-pipeline.spec.json}{Rendering Pipeline Spec}
\end{itemize}

\section{Shells, DOGFOOD, And Proving Apps}\label{record-001-shells-dogfood-and-proving-apps}

DOGFOOD is not a side example. It is the canonical human-facing docs and shell proving surface for the repo.

Use these docs when you are shaping shell behavior, shell-owned overlays, settings, notification review, or documentation posture:

\begin{itemize}
\tightlist
\item
  \href{https://github.com/flyingrobots/bijou/blob/cbd70612140e7a450d764b3c7be1c6af878b99bc/docs/DOGFOOD.md}{DOGFOOD}
\item
  \href{https://github.com/flyingrobots/bijou/blob/cbd70612140e7a450d764b3c7be1c6af878b99bc/docs/strategy/dogfood.md}{DOGFOOD strategy note}
\item
  \href{https://github.com/flyingrobots/bijou/blob/cbd70612140e7a450d764b3c7be1c6af878b99bc/docs/strategy/humane-shell.md}{Humane Shell}
\item
  \href{https://github.com/flyingrobots/bijou/blob/cbd70612140e7a450d764b3c7be1c6af878b99bc/docs/strategy/settings-belong-to-the-shell.md}{Settings Belong To The Shell}
\item
  \href{https://github.com/flyingrobots/bijou/blob/cbd70612140e7a450d764b3c7be1c6af878b99bc/docs/strategy/notification-history-belongs-to-the-shell.md}{Notification History Belongs To The Shell}
\item
  \href{https://github.com/flyingrobots/bijou/blob/cbd70612140e7a450d764b3c7be1c6af878b99bc/packages/bijou-tui/ADVANCED_GUIDE.md}{packages/bijou-tui/ADVANCED\_GUIDE.md}
\end{itemize}

Useful proving surfaces:

\begin{itemize}
\tightlist
\item
  \href{https://github.com/flyingrobots/bijou/blob/cbd70612140e7a450d764b3c7be1c6af878b99bc/examples/docs/app.ts}{examples/docs/app.ts}
\item
  \href{https://github.com/flyingrobots/bijou/blob/cbd70612140e7a450d764b3c7be1c6af878b99bc/examples/release-workbench/main.ts}{examples/release-workbench/main.ts}
\item
  \href{https://github.com/flyingrobots/bijou/blob/cbd70612140e7a450d764b3c7be1c6af878b99bc/scripts/docs-preview.test.ts}{scripts/docs-preview.test.ts}
\item
  \href{https://github.com/flyingrobots/bijou/blob/cbd70612140e7a450d764b3c7be1c6af878b99bc/tests/cycles/RE-017/frame-shell-theme-dogfood-demo.test.ts}{tests/cycles/RE-017/frame-shell-theme-dogfood-demo.test.ts}
\end{itemize}

\section{Pipeline, Shaders, Painting, And Transforms}\label{record-001-pipeline-shaders-painting-and-transforms}

If you are touching frame-to-frame rendering, post-process effects, or shader behavior, do not stop at the API entrypoints. Read the doctrine and the implementation seams together.

Start here:

\begin{itemize}
\tightlist
\item
  \href{https://github.com/flyingrobots/bijou/blob/cbd70612140e7a450d764b3c7be1c6af878b99bc/packages/bijou-tui/ADVANCED_GUIDE.md}{packages/bijou-tui/ADVANCED\_GUIDE.md}
\item
  \href{https://github.com/flyingrobots/bijou/blob/cbd70612140e7a450d764b3c7be1c6af878b99bc/docs/strategy/layout-localization-pipeline.md}{Layout Localization Pipeline}
\item
  \href{https://github.com/flyingrobots/bijou/blob/cbd70612140e7a450d764b3c7be1c6af878b99bc/packages/bijou-tui/src/pipeline/pipeline.ts}{packages/bijou-tui/src/pipeline/pipeline.ts}
\item
  \href{https://github.com/flyingrobots/bijou/blob/cbd70612140e7a450d764b3c7be1c6af878b99bc/packages/bijou-tui/src/pipeline/middleware/paint.ts}{paint middleware}
\item
  \href{https://github.com/flyingrobots/bijou/blob/cbd70612140e7a450d764b3c7be1c6af878b99bc/packages/bijou-tui/src/pipeline/middleware/motion.ts}{motion middleware}
\item
  \href{https://github.com/flyingrobots/bijou/blob/cbd70612140e7a450d764b3c7be1c6af878b99bc/packages/bijou-tui/src/pipeline/middleware/css.ts}{css middleware}
\item
  \href{https://github.com/flyingrobots/bijou/blob/cbd70612140e7a450d764b3c7be1c6af878b99bc/packages/bijou-tui/src/pipeline/middleware/grayscale.ts}{grayscale middleware}
\item
  \href{https://github.com/flyingrobots/bijou/blob/cbd70612140e7a450d764b3c7be1c6af878b99bc/packages/bijou-tui/src/transition-shaders.ts}{transition shaders}
\item
  \href{https://github.com/flyingrobots/bijou/blob/cbd70612140e7a450d764b3c7be1c6af878b99bc/docs/specs/shader-v2.spec.json}{Shader V2 Spec}
\item
  \href{https://github.com/flyingrobots/bijou/blob/cbd70612140e7a450d764b3c7be1c6af878b99bc/docs/BACKLOG/inbox/DX-008-render-pipeline-guide.md}{Render Pipeline guide backlog note}
\end{itemize}

\section{Performance And Soak Work}\label{record-001-performance-and-soak-work}

Bijou treats performance as a product surface, not as a cleanup pass.

The current performance and soak lane is spread across:

\begin{itemize}
\tightlist
\item
  \href{https://github.com/flyingrobots/bijou/blob/cbd70612140e7a450d764b3c7be1c6af878b99bc/examples/perf-gradient/main.ts}{examples/perf-gradient/main.ts}
\item
  \href{https://github.com/flyingrobots/bijou/blob/cbd70612140e7a450d764b3c7be1c6af878b99bc/packages/bijou/src/core/components/perf-overlay.ts}{perf overlay}
\item
  \href{https://github.com/flyingrobots/bijou/blob/cbd70612140e7a450d764b3c7be1c6af878b99bc/docs/perf/RE-017-byte-pipeline.md}{Byte-Pipeline Recovery}
\item
  \href{https://github.com/flyingrobots/bijou/blob/cbd70612140e7a450d764b3c7be1c6af878b99bc/docs/releases/4.4.0/whats-new.md}{4.4.0 release notes}
\item
  \href{https://github.com/flyingrobots/bijou/blob/cbd70612140e7a450d764b3c7be1c6af878b99bc/docs/BACKLOG/cool-ideas/DX-006-debug-performance-overlay-component.md}{Debug Performance Overlay Component}
\end{itemize}

That is the lane to read when you are investigating hot-loop cost, proving a new optimization, or deciding whether a feature belongs in the runtime, middleware, or shell.

\section{Localization And I18n Workflows}\label{record-001-localization-and-i18n-workflows}

Localization is not only string substitution. The repo posture includes directionality, shell copy, proving surfaces, and exchange workflows.

Start here:

\begin{itemize}
\tightlist
\item
  \href{https://github.com/flyingrobots/bijou/blob/cbd70612140e7a450d764b3c7be1c6af878b99bc/docs/strategy/localization-and-bidirectionality.md}{Localization and Bidirectionality}
\item
  \href{https://github.com/flyingrobots/bijou/blob/cbd70612140e7a450d764b3c7be1c6af878b99bc/docs/strategy/layout-localization-pipeline.md}{Layout Localization Pipeline}
\item
  \href{https://github.com/flyingrobots/bijou/blob/cbd70612140e7a450d764b3c7be1c6af878b99bc/packages/bijou-i18n/README.md}{packages/bijou-i18n/README.md}
\item
  \href{https://github.com/flyingrobots/bijou/blob/cbd70612140e7a450d764b3c7be1c6af878b99bc/packages/bijou-i18n-tools/README.md}{packages/bijou-i18n-tools/README.md}
\item
  \href{https://github.com/flyingrobots/bijou/blob/cbd70612140e7a450d764b3c7be1c6af878b99bc/packages/bijou-i18n-tools-node/README.md}{packages/bijou-i18n-tools-node/README.md}
\item
  \href{https://github.com/flyingrobots/bijou/blob/cbd70612140e7a450d764b3c7be1c6af878b99bc/packages/bijou-i18n-tools-xlsx/README.md}{packages/bijou-i18n-tools-xlsx/README.md}
\item
  \href{https://github.com/flyingrobots/bijou/blob/cbd70612140e7a450d764b3c7be1c6af878b99bc/docs/design/LX-001-bijou-i18n-runtime-package.md}{LX-001 runtime package design}
\item
  \href{https://github.com/flyingrobots/bijou/blob/cbd70612140e7a450d764b3c7be1c6af878b99bc/docs/design/LX-002-bijou-i18n-tools-package.md}{LX-002 tooling package design}
\end{itemize}

When validating localization work, also look at the proving surfaces:

\begin{itemize}
\tightlist
\item
  \href{https://github.com/flyingrobots/bijou/blob/cbd70612140e7a450d764b3c7be1c6af878b99bc/examples/docs/app.ts}{examples/docs/app.ts}
\item
  \href{https://github.com/flyingrobots/bijou/blob/cbd70612140e7a450d764b3c7be1c6af878b99bc/scripts/docs-preview.test.ts}{scripts/docs-preview.test.ts}
\item
  \href{https://github.com/flyingrobots/bijou/blob/cbd70612140e7a450d764b3c7be1c6af878b99bc/tests/cycles/LX-008/localized-shell-and-dogfood.test.ts}{tests/cycles/LX-008/localized-shell-and-dogfood.test.ts}
\end{itemize}

\section{Runtime Boundaries, Workers, And Capture}\label{record-001-runtime-boundaries-workers-and-capture}

If the question is about Node ownership, worker threads, recorder flows, or how Bijou crosses the runtime boundary, read the Node lane directly:

\begin{itemize}
\tightlist
\item
  \href{https://github.com/flyingrobots/bijou/blob/cbd70612140e7a450d764b3c7be1c6af878b99bc/packages/bijou-node/ADVANCED_GUIDE.md}{packages/bijou-node/ADVANCED\_GUIDE.md}
\item
  \href{https://github.com/flyingrobots/bijou/blob/cbd70612140e7a450d764b3c7be1c6af878b99bc/docs/specs/worker-runtime.spec.json}{worker runtime spec}
\item
  \href{https://github.com/flyingrobots/bijou/blob/cbd70612140e7a450d764b3c7be1c6af878b99bc/examples/v3-worker/README.md}{examples/v3-worker/README.md}
\item
  \href{https://github.com/flyingrobots/bijou/blob/cbd70612140e7a450d764b3c7be1c6af878b99bc/examples/v3-worker/main.ts}{examples/v3-worker/main.ts}
\item
  \href{https://github.com/flyingrobots/bijou/blob/cbd70612140e7a450d764b3c7be1c6af878b99bc/examples/v3-worker/record.ts}{examples/v3-worker/record.ts}
\end{itemize}

\section{Testing, PTY, And Release Proof}\label{record-001-testing-pty-and-release-proof}

The advanced bar in Bijou is not ``unit tests only.'' Real proof usually spans:

\begin{itemize}
\tightlist
\item
  deterministic package tests
\item
  DOGFOOD/docs preview tests
\item
  PTY or canary smoke coverage
\item
  release-readiness and release-preflight
\item
  targeted performance probes when render-path work changes
\end{itemize}

Useful anchors:

\begin{itemize}
\tightlist
\item
  \href{https://github.com/flyingrobots/bijou/blob/cbd70612140e7a450d764b3c7be1c6af878b99bc/scripts/smoke-canaries.ts}{scripts/smoke-canaries.ts}
\item
  \href{https://github.com/flyingrobots/bijou/blob/cbd70612140e7a450d764b3c7be1c6af878b99bc/scripts/smoke-dogfood.test.ts}{scripts/smoke-dogfood.test.ts}
\item
  \href{https://github.com/flyingrobots/bijou/blob/cbd70612140e7a450d764b3c7be1c6af878b99bc/scripts/docs-preview.test.ts}{scripts/docs-preview.test.ts}
\item
  \href{https://github.com/flyingrobots/bijou/blob/cbd70612140e7a450d764b3c7be1c6af878b99bc/docs/release.md}{docs/release.md}
\item
  \href{https://github.com/flyingrobots/bijou/blob/1853ad902f6bf219326409e7f743a5a1d6517dbc/docs/method/release-runbook.md}{docs/method/release-runbook.md}
\end{itemize}

\section{Rule Of Thumb}\label{record-001-rule-of-thumb}

If the public \texttt{G\allowbreak{}U\allowbreak{}I\allowbreak{}D\allowbreak{}E\allowbreak{}.\allowbreak{}m\allowbreak{}d} got you productive, stay there.

If you are asking ``how does the shell own this,'' ``where should this render-path decision live,'' ``how do we prove this in DOGFOOD,'' or ``what workflow keeps this truthful across runtime modes,'' you are in \texttt{A\allowbreak{}D\allowbreak{}V\allowbreak{}A\allowbreak{}N\allowbreak{}C\allowbreak{}E\allowbreak{}D\allowbreak{}\_\allowbreak{}G\allowbreak{}U\allowbreak{}I\allowbreak{}D\allowbreak{}E} territory.

